% =============================================================================
%  Chapter 20 — VTK Export Cleanup
% =============================================================================

\chapter{VTK Export Cleanup}
\label{ch:vtk_cleanup}

This chapter documents the cleanup pass that reduced the VTK layer to the
smallest operational contract that still preserves extension points.


% ─────────────────────────────────────────────────────────────────────────
\section{Retained Responsibilities}
\label{sec:vtk_cleanup_retained}

After the cleanup, the active VTK surface in \texttt{fall\_n} is:
\begin{itemize}[nosep]
  \item \texttt{VTKCellTraits.hh}: mapping from \texttt{fall\_n} topology to
        VTK cell ids and local-node permutations;
  \item \texttt{VTKTensorFieldDerivatives.hh}: deterministic tensor-to-scalar
        post-processing utilities;
  \item \texttt{VTKModelExporter.hh}: split export of the continuum mesh and
        the Gauss cloud;
  \item \texttt{PVDWriter.hh}: time-series collection writer.
\end{itemize}

This is enough for the current library needs:
\begin{itemize}[nosep]
  \item clean nodal contour plots on the mesh;
  \item faithful material-point inspection on a warpable Gauss cloud;
  \item explicit scalar invariants for ParaView;
  \item time-series support through \texttt{.pvd}.
\end{itemize}


% ─────────────────────────────────────────────────────────────────────────
\section{Removed Complexity}
\label{sec:vtk_cleanup_removed}

The cleanup explicitly removed three categories of dead or unstable
implementation:
\begin{enumerate}[nosep]
  \item embedded native quadrature export as part of the main exporter API;
  \item the mixed single-file visualization mesh that combined continuum cells
        and Gauss glyphs in one \texttt{.vtu};
  \item broad VTK rendering and GUI dependencies from the common include path.
\end{enumerate}

These paths were not removed because they were impossible.  They were removed
because they increased API surface, state management, and cognitive load
without yielding a robust everyday workflow in ParaView.


% ─────────────────────────────────────────────────────────────────────────
\section{Why the Split Export Is the Right Current Boundary}
\label{sec:vtk_cleanup_boundary}

The split contract is not merely a fallback; it is the correct current
boundary between FE post-processing and VTK:
\begin{itemize}[nosep]
  \item the continuum mesh remains semantically clean;
  \item the Gauss cloud carries only material-point quantities;
  \item \texttt{Warp By Vector} works directly on both datasets;
  \item ParaView is not asked to reconstruct semantics from ambiguous metadata;
  \item the exporter state is easy to test and reason about.
\end{itemize}

This matters architecturally.  A post-processing layer should not be clever
just because VTK allows it; it should be explicit enough that future extensions
do not inherit accidental complexity.


% ─────────────────────────────────────────────────────────────────────────
\section{Compile-Time and Dependency Impact}
\label{sec:vtk_cleanup_compiletime}

The cleanup also reduced compilation and linkage overhead:
\begin{itemize}[nosep]
  \item \texttt{header\_files.hh} no longer pulls the broad
        \texttt{VTKheaders.hh} compatibility umbrella;
  \item the compatibility umbrella itself was reduced to minimal data-model and
        IO includes for legacy code;
  \item the CMake VTK dependency set now targets only the modules needed by the
        active exporter and tests.
\end{itemize}

This follows the library's general HPC-oriented guideline: keep heavyweight and
GUI-oriented dependencies out of the transitive include path unless they are
strictly required.


% ─────────────────────────────────────────────────────────────────────────
\section{Retained Extension Points}
\label{sec:vtk_cleanup_extensions}

The cleanup did \emph{not} freeze the VTK layer.  The intended extension points
remain:
\begin{itemize}[nosep]
  \item adding new explicit scalar filters from tensor data;
  \item adding new Gauss-cloud fields as constitutive models evolve;
  \item extending \texttt{VTKCellTraits} as new geometries are introduced;
  \item adding alternate writers as isolated adapters rather than as branches
        inside the main exporter.
\end{itemize}

That is the key design rule going forward: extend through small orthogonal
adapters, not by re-expanding \texttt{VTKModelExporter} into a multipurpose
state machine.

